\documentclass[a4paper,11pt]{article}
\usepackage[utf8]{inputenc}
\usepackage[T1]{fontenc}
\usepackage[french]{babel}
\usepackage{graphicx}
\usepackage{amsmath}
\usepackage{amsthm}
\theoremstyle{definition}
\usepackage{geometry}
\usepackage{fancyhdr}
\usepackage{hyperref}
\usepackage[usenames,dvipsnames,svgnames,table]{xcolor}
\usepackage{enumerate}
\usepackage{tcolorbox}
\tcbuselibrary{skins,breakable}
\usepackage{tikz}
\usepackage{pgfplots}
\pgfplotsset{compat=1.18}
\usetikzlibrary{arrows.meta}

\geometry{a4paper, scale=0.8}

%% ─── Header / footer ─────────────────────────────────────────────
\pagestyle{fancy}
\fancyhf{}
\fancyhead[L]{2025-2026}
\fancyhead[C]{\textbf{\courseName}}
\fancyhead[R]{\large GEII -- BUT2}
\fancyfoot[C]{\thepage}
\fancyfoot[L]{\includegraphics[height=3em]{../../common/logos_couleurs/IUTA.png}}
% \fancyfoot[R]{\href{mailto:prenom.nom@univ-smb.fr}{prenom.nom@univ-smb.fr}}

%% ─── Course metadata ─────────────────────────────────────────────
\newcommand{\courseName}{Outils Mathématiques et Logiciel --- OML3}
\newcommand{\chapterTitle}{TD 2 -- Série de Fourier}

%% ─── Custom title ────────────────────────────────────────────────
\makeatletter
\renewcommand\maketitle{{\raggedright\begin{center}
    {\Huge\bfseries\@title}
\end{center}}}
\makeatother

%% ─── Solution toggle ─────────────────────────────────────────────
%% Compile with \showsolutionstrue for the teacher version.
\newif\ifshowsolutions
\showsolutionsfalse
\ifshowsolutions
    \newtcolorbox{solution}{
        colback=white, colframe=white, coltext=red,
        breakable, title=\textbf{Solution :}
    }
\else
    \NewEnviron{solution}{}{}
\fi

%% ─── Theorem / exercise environments ─────────────────────────────
\newtheorem{definition}{Définition}[section]
\newtheorem{proposition}{Propriété}[section]
\newtheorem{theorem}{Théorème}[section]
\newtheorem{exemple}{Exemple}[section]
\newtheorem{remark}{Remarque}[section]

%% Exercice : affiche "Exercice N (titre)." et referme aussitôt le
%% paragraphe avec \par. Volontairement PAS un \newtheorem : amsthm
%% reporte l'affichage du titre via \everypar, qui ne se déclenche
%% qu'au démarrage d'un paragraphe. Si le corps commence directement
%% par une liste (\begin{enumerate}), aucun paragraphe ne démarre
%% avant le premier \item, et le titre se retrouve collé au "1." de
%% la liste -- même avec une ligne vide dans la source. Cet
%% environnement affiche le titre immédiatement, donc ne dépend pas
%% de ce mécanisme et fonctionne quel que soit le contenu qui suit.
\newcounter{exercice}   % global counter — do not add [section]
\newenvironment{exercice}[1][]{%
  \refstepcounter{exercice}%
  \par\noindent\textbf{Exercice \theexercice}\if\relax\detokenize{#1}\relax\else\ (#1)\fi.\par\noindent\ignorespaces
}{\par}

%% ─── Section numbering ───────────────────────────────────────────
\renewcommand{\thesection}{{\arabic{section}}}

%% ─── Helpers ─────────────────────────────────────────────────────
\definecolor{niceblue}{RGB}{0,128,172}
\newcommand{\highlight}[1]{{\color{niceblue}\emph{#1}}}

\usepackage{amsmath,amsfonts,bm}

%% ─── Operators ───────────────────────────────────────────────────
\DeclareMathOperator*{\argmax}{arg\,max}
\DeclareMathOperator*{\argmin}{arg\,min}
\DeclareMathOperator{\MSE}{MSE}
\DeclareMathOperator{\sign}{sign}
\DeclareMathOperator{\sinc}{sinc}
\DeclareMathOperator{\tr}{tr}

%% ─── Number sets ─────────────────────────────────────────────────
\def\R{{\mathbb{R}}}
\def\C{{\mathbb{C}}}
\def\N{{\mathbb{N}}}
\def\Z{{\mathbb{Z}}}

%% ─── Infinitesimal and units ────────────────────────────────────
\def\d{{\rm d}}
\def\Hz{{\mbox{Hz}}}
\def\m{{\mbox{m}}}

%% ─── Complex / signal ───────────────────────────────────────────
\DeclareMathOperator{\complexj}{\mathsf{j}}   % imaginary unit

%% ─── Probability ─────────────────────────────────────────────────
\newcommand{\Esp}[2]{\mathbb{E}_{#1}\!\left[#2\right]}
\renewcommand{\Pr}{\mathbb{P}}

%% ─── Linear algebra ──────────────────────────────────────────────
\newcommand{\abs}[1]{\left|{#1}\right|}
\newcommand{\norm}[1]{\left\|#1\right\|}
\newcommand{\inner}[2]{\left\langle#1,\,#2\right\rangle}

%% ─── Logic ───────────────────────────────────────────────────────
\newcommand{\then}{\Rightarrow}
\newcommand{\iif}{\Leftrightarrow}

%% ─── Draft annotation ────────────────────────────────────────────
\newcommand{\red}[1]{{\color{red}#1}}

%% ─── Shared figure color (used by common/figures/*/*.tex bodies) ──
\definecolor{niceblue}{RGB}{0,128,172}

%% ─── Rappel header (styled subtitle for reference-card asides) ───
%% Requires \definecolor{niceblue}{...} to be defined by the including document.
\newcommand{\rappelhead}[1]{%
    \bigskip
    {\color{niceblue}\bfseries #1}\\[-0.4em]
    {\color{niceblue!35}\rule{\linewidth}{0.6pt}}
    \smallskip

}


\title{\chapterTitle}
\author{}
\date{}

%% ─────────────────────────────────────────────────────────────────
\begin{document}
\fontfamily{bch}\selectfont
\fontsize{12pt}{14pt}\selectfont

\maketitle
\thispagestyle{fancy}
\setcounter{exercice}{0}


\vspace{2em}
\textit{Tous les exercices ne seront pas abordés lors des séances de TD. Pour vos révisions il est recommandé de revoir précisément les 3 premiers exercices de chaque TD et de vous confronter au moins une fois à chaque exercice. La correction des TD sera mise en ligne sur le Moodle du cours dans la semaine qui suit le TD.}
\vspace{2em}


\begin{exercice}[Décalage dans le temps]
    Soit $f$ une fonction $T$-périodique, de coefficients de Fourier complexes $c_n$. Pour $\varphi\in\R$, on considère le signal décalé $g(t) = f(t+\varphi)$.

    \begin{enumerate}[\bfseries 1.]
        \item Montrer que $g$ est $T$-périodique.
        \begin{solution}
            $g(t+T) = f(t+T+\varphi) = f(t+\varphi) = g(t)$ car $f$ est $T$-périodique : $g$ est donc $T$-périodique, de même période que $f$.
        \end{solution}

    

        \item En notant $d_n$ les coefficients de Fourier complexes de $g$, montrer que $d_n = c_n\,e^{jn\omega_0\varphi}$.
        \begin{solution}
            $$
                d_n = \frac1T\int_0^T g(t)e^{-jn\omega_0t}\,\d t = \frac1T\int_0^T f(t+\varphi)e^{-jn\omega_0t}\,\d t.
            $$
            On pose $u=t+\varphi$ (donc $\d u=\d t$) :
            $$
                d_n = \frac1T\int_\varphi^{T+\varphi} f(u)e^{-jn\omega_0(u-\varphi)}\,\d u = e^{jn\omega_0\varphi}\cdot\frac1T\int_\varphi^{T+\varphi} f(u)e^{-jn\omega_0u}\,\d u.
            $$
            La fonction $u\mapsto f(u)e^{-jn\omega_0u}$ est $T$-périodique (produit de $f$, $T$-périodique, et de $u\mapsto e^{-jn\omega_0u}$, $T$-périodique puisque $\omega_0T=2\pi$). Son intégrale sur $[\varphi,T+\varphi]$ vaut la même chose que sur $[0,T]$ (leur donner ce résultat s'il ne leur paraît pas évident puis les rediriger vers l'exercice 5 du TD1 pour la démo), donc
            $$
                d_n = e^{jn\omega_0\varphi}\cdot\frac1T\int_0^T f(u)e^{-jn\omega_0u}\,\d u = c_n\,e^{jn\omega_0\varphi}.
            $$
        \end{solution}

        % \item En utilisant la formule d'Euler, calculer les coefficients $c_n$ de $h(t)=\cos(\omega_0t)$ et tracer son spectre bilatéral.
        % \begin{solution}
        %     $\cos(\omega_0t) = \dfrac{e^{j\omega_0t}+e^{-j\omega_0t}}{2}$ (Propriété~2.6). Par identification avec $\sum_n c_n e^{jn\omega_0t}$,
        %     $$
        %         c_1 = c_{-1} = \frac12, \qquad c_n = 0 \text{ pour } n\notin\{-1,1\}.
        %     $$
        %     Le spectre bilatéral ne comporte donc que deux bâtons, de hauteur $|c_1|=|c_{-1}|=\frac12$, aux fréquences $\pm f_0$ :
        %     \begin{center}
        %     \begin{tikzpicture}
        %         \begin{axis}[
        %             width=7cm, height=3.5cm,
        %             xlabel={fréquence}, ylabel={$|c_n|$},
        %             xtick={-1,0,1}, xticklabels={$-f_0$,$0$,$f_0$},
        %             ytick={0.5}, yticklabels={$\frac12$},
        %             ymin=0, ymax=0.7,
        %             axis lines=left,
        %         ]
        %         \addplot+[ycomb, niceblue, thick, mark=*, mark options={fill=niceblue}] coordinates {(-1,0.5) (0,0) (1,0.5)};
        %         \end{axis}
        %     \end{tikzpicture}
        %     \end{center}

            \item Calculer les coefficients complexes de la fonction carré définie sur une période $T>0$ par
        $$
            f_{carre}(t) = \begin{cases}
                1 & t\in[0,T/2[\\ -1 & t\in[T/2, T[
                    \end{cases}
        $$ 

        \begin{solution}
            Soit $n\in\Z$
            $$
                c_n = \frac{1}{T}\int_0^T f_{carre}(t) e^{-\complexj n \omega t} \d t = \frac{1}{T}\int_0^{T/2} e^{-\complexj n \omega t} \d t - \frac{1}{T}\int_{T/2}^T e^{-\complexj n \omega t} \d t.
            $$
            On applique le changement de variable $u= t-T/2$ de sorte à avoir
            $$
            \begin{aligned}
                c_n &= \frac{1}{T}\int_0^{T/2} e^{-\complexj n \omega t} \d t - \frac{1}{T}\int_{0}^{T/2} e^{-\complexj n \omega (t + T/2)} \d t \\
                &= \frac{1}{T}\int_0^{T/2} e^{-\complexj n \omega t} \d t - \frac{e^{-\complexj n \omega \pi}}{T} \int_{0}^{T/2} e^{-\complexj n \omega t}  \d t
            \end{aligned}
            $$
            ($\omega T/2 = \pi$).
            On distingue alors deux cas. 

            \textbf{$n$ paire :} $e^{-\complexj n \pi} = 1$ et donc $c_n = 0$.
            
            \textbf{$n$ impaire :} $e^{-\complexj n \pi} = -1$ et donc 
            $$
            \begin{aligned}
                c_n &= \frac{2}{T} \int_0^{T/2} e^{-\complexj n \omega t} \d t = \frac{2}{T} \left[ \frac{e^{-\complexj n\omega t}}{-\complexj n \omega}\right]_0^{T/2} \\
                &= \frac{2}{T} \left( \frac{-1 - 1}{-\complexj n\omega}\right) = -\complexj \frac{2}{n \pi}.
            \end{aligned}
            $$
            (En bonus pour ceux qui s'ennuie : en déduire les coefs $a_n$ et $b_n$ de la forme trigonométrique -> identifier à partir de $c_n = \frac{a_n - j b_n}{2}$).

        \end{solution}

        \item En déduire l'expression des coefficients $c_n$ de la fonction périodique définie par 
        $$
        g(t) = 3f_{\rm carre}(t) + f_{\rm carre}(t-2)
        $$
        et tracer son spectre bilatéral pour $-8 \leq n \leq 8$.
        \begin{solution}
            En appliquant la formule précédemment démontré et en utlisant la linéarité de la série de Fourier : 
            $$
            \begin{aligned}    
                d_n &= \begin{cases}
                    0 & n \ \text{paire} \\
                    -3\complexj \frac{2}{n \pi} - \complexj e^{\complexj 2 n\omega}  \frac{2}{n \pi}  & n \ \text{impaire} \\
                \end{cases}\\ 
                d_n &= \begin{cases}
                    0 & n \ \text{paire} \\
                    \frac{2}{n \pi}  \left( 3 e^{-\pi/2} + e^{\complexj 2 n\omega - \pi/2}\right)  & n \ \text{impaire} \\
                \end{cases}
            \end{aligned}
                $$

            (En bonus : Est-ce $d_n$ imaginaire pure ? -> non. Que conclure sur $g$ ? -> n'est pas paire). 
        \end{solution}

        % \item En déduire les coefficients de $\sin(\omega_0t)$ à l'aide de la question 2 puis tracer son spectre bilatéral.
        % \begin{solution}
        %     $\sin(\omega_0t) = \cos\!\left(\omega_0t-\dfrac\pi2\right) = h\!\left(t+\varphi\right)$ avec $\varphi = -\dfrac{\pi}{2\omega_0}$. D'après la question 2, les coefficients de $\sin(\omega_0t)$ sont $d_n=c_n e^{jn\omega_0\varphi}$ :
        %     $$
        %         d_1 = \frac12 e^{j\omega_0\times(-\pi/(2\omega_0))} = \frac12 e^{-j\pi/2} = -\frac{j}2, \qquad d_{-1} = \frac12 e^{j\pi/2} = \frac{j}2, \qquad d_n=0 \text{ sinon.}
        %     $$
        %     On retrouve bien, via la formule d'Euler, $\sin(\omega_0t) = \dfrac{e^{j\omega_0t}-e^{-j\omega_0t}}{2j} = -\dfrac{j}{2}e^{j\omega_0t} + \dfrac{j}{2}e^{-j\omega_0t}$.

        %     Ici $|d_1|=|d_{-1}|=\frac12$ : le spectre bilatéral de $\sin(\omega_0t)$ est donc \emph{identique} à celui de $\cos(\omega_0t)$ (deux bâtons de hauteur $\frac12$ en $\pm f_0$). Le spectre d'amplitude ne conserve que le module des $c_n$ : il ne distingue pas $\cos$ et $\sin$, seule leur phase (l'argument de $d_n$) diffère.
        % \end{solution}
    \end{enumerate}
\end{exercice}
%% ── Ex 1 : PROVE (guided) ────────────────────────────────────────
%% Purpose: establish the one lemma every other exercise in this TD
%% leans on (isolating a term via orthogonality of complex
%% exponentials over a period), guided step by step so students derive
%% it themselves rather than being handed it.
\begin{exercice}
    Soit $T>0$, $\omega_0=\dfrac{2\pi}{T}$, et $k\in\Z$. On veut établir le résultat suivant
    $$
        \int_0^T e^{\complexj k\omega_0 t}\,\d t = \begin{cases} T & \text{si } k=0 \\ 0 & \text{si } k\neq 0 \end{cases}.
    $$

    \begin{enumerate}[\bfseries 1.]
        \item Traiter le cas $k=0$.
        \begin{solution}
            Pour $k=0$, $e^{\complexj k\omega_0 t}=1$, donc $\displaystyle\int_0^T e^{\complexj k\omega_0 t}\,\d t = \int_0^T 1\,\d t = T$.
        \end{solution}

        \item On suppose maintenant $k\neq 0$. Donner une primitive de $t\mapsto e^{\complexj k\omega_0 t}$.
        \begin{solution}
            $t\mapsto \dfrac{e^{\complexj k\omega_0 t}}{\complexj k\omega_0}$ est une primitive de $t\mapsto e^{\complexj k\omega_0 t}$ (on dérive comme une exponentielle usuelle ; $k\omega_0\neq0$ assure que l'on ne divise pas par $0$).
        \end{solution}

        \item En déduire $\displaystyle\int_0^T e^{\complexj k\omega_0 t}\,\d t$ pour $k\neq0$. On pourra utiliser $\omega_0 T = 2\pi$.
        \begin{solution}
            $$
                \int_0^T e^{\complexj k\omega_0 t}\,\d t = \left[\frac{e^{\complexj k\omega_0 t}}{\complexj k\omega_0}\right]_0^T = \frac{e^{\complexj k\omega_0 T}-1}{\complexj k\omega_0}.
            $$
            Or $k\omega_0T = 2k\pi$ avec $k\in\Z$, donc $e^{\complexj k\omega_0 T} = e^{\complexj 2k\pi} = 1$, d'où
            $$
                \int_0^T e^{\complexj k\omega_0 t}\,\d t = \frac{1-1}{\complexj k\omega_0} = 0.
            $$
        \end{solution}

        % \item En séparant partie réelle et partie imaginaire, retrouver que $\displaystyle\int_0^T \cos(k\omega_0 t)\,\d t = 0$ et $\displaystyle\int_0^T \sin(k\omega_0 t)\,\d t = 0$ pour tout $k\in\Z^*$.
        % \begin{solution}
        %     $e^{\complexj k\omega_0 t} = \cos(k\omega_0 t) + \complexj\sin(k\omega_0 t)$ (formule d'Euler), donc
        %     $$
        %         \int_0^T e^{\complexj k\omega_0 t}\,\d t = \int_0^T \cos(k\omega_0 t)\,\d t + \complexj\int_0^T \sin(k\omega_0 t)\,\d t.
        %     $$
        %     Cette quantité est nulle (question précédente) ; une somme $a+\complexj b$ ($a,b\in\R$) n'est nulle que si $a=b=0$. D'où $\displaystyle\int_0^T \cos(k\omega_0 t)\,\d t = 0$ et $\displaystyle\int_0^T \sin(k\omega_0 t)\,\d t = 0$.
        % \end{solution}
    \end{enumerate}

\end{exercice}


%% ── Ex 2 : CLASSIFY / DEFINE ─────────────────────────────────────
%% Purpose: activate vocabulary from the CM.
%% No computation harder than substituting into a definition.
%% At least one positive, one negative, one boundary case.
\begin{exercice}
        On considère la fonction $$f(t) = \cos(\omega_0 t + \varphi),$$ où $\omega_0$ et $\varphi$ sont deux réels.

        \begin{enumerate}[\bfseries 1.]
            \item $f$ est-elle périodique et si oui ecrire sa période en fonction de $\omega_0$.
             \begin{solution}
                On vérifie que $f$ est périodique de période $T = \frac{2\pi}{\omega_0}$.
             \end{solution}
             \item Montrer que $f$ est paire si et seulement si $\varphi = k\pi$, $k\in\Z$
            \begin{solution}
             $f$ est paire si et seulement si, pour tout réel $t$
             $
             f(t) = f(-t),
             $ soit
             $$
             \begin{aligned}
                &\cos(\omega_0 t + \varphi) = \cos(-\omega_0 t + \varphi) \\
                \iff & \cos(\omega_0 t)\cos(\varphi) - \sin(\omega_0 t) \sin(\varphi) = \cos(\omega_0t) \cos(\varphi) + \sin(\omega_0 t)\sin(\varphi) \\
                \iff & \sin(\varphi) = 0.
             \end{aligned}
             $$
             Ce qui équivaut à $\varphi = k\pi$, $k\in\Z$.
             \end{solution}


            \item Montrer que $f$ est impaire si et seulement si $\varphi = \frac{k\pi}{2}$ , $k\in\Z$.
            \begin{solution}
                De la même façon,  $f$ est impaire si et seulement si, pour tout réel $t$
             $
             f(t) = -f(-t),
             $ soit
             $$
             \begin{aligned}
                &\cos(\omega_0 t + \varphi) = -\cos(-\omega_0 t + \varphi) \\
                \iff & \cos(\omega_0 t)\cos(\varphi) - \sin(\omega_0 t) \sin(\varphi) = -\cos(\omega_0t) \cos(\varphi) - \sin(\omega_0 t)\sin(\varphi) \\
                \iff & \cos(\varphi) = 0.
             \end{aligned}
             $$
             Ce qui est équivaut à $\varphi = k\cfrac{\pi}{2}$, $k\in\Z$.
            \end{solution}

            \item Que valent les coefficients complexes de la série de Fourier de $f$ pour $n \neq 1$ et $n\neq -1$.
            \begin{solution}
                On rappel la formule des coefficients complexes. Pour $n\in\Z$
                $$
                    c_n = \frac{1}{T} \int_0^T f(t) e^{-jn\omega_0t}\d t.
                $$
                Plutôt que de chercher la primitive du produit, on exprimera la fonction cosinus par la formule d'Euler
                $$
                    \cos(t) = \frac{e^{\complexj t} + e^{-\complexj t}}{2}.
                $$
                On obtient
                $$
                \begin{aligned}
                    c_n &= \frac{1}{T}\int_0^T \left( \frac{e^{\complexj \omega_0 t + \varphi} + e^{-\complexj \omega_0 t + \varphi}}{2} \right) e^{-jn\omega_0t}\d t = \frac{e^{\complexj\varphi}}{T}\int_0^T \left( \frac{e^{\complexj \omega_0 t} + e^{-\complexj \omega_0 t}}{2} \right) e^{-jn\omega_0t}\d t \\
                    & = \frac{e^{\complexj\varphi}}{2T}\int_0^T e^{-\complexj (n-1)\omega_0 t} + e^{-\complexj (n+1)\omega_0 t} \d t = \frac{e^{\complexj\varphi}}{2T} \left(\left[- \frac{e^{-\complexj (n-1)\omega_0 t}}{\complexj (n-1)\omega_0}\right]_0^T + \left[- \frac{e^{-\complexj (n+1)\omega_0 t}}{\complexj (n+1)\omega_0}\right]_0^T\right) \\
                    &= \frac{e^{\complexj\varphi}}{2T}\left(\frac{1 - e^{-\complexj (n-1)\omega_0 T}}{\complexj (n-1)\omega_0} + \frac{1 - e^{-\complexj (n+1)\omega_0 T}}{\complexj (n+1)\omega_0}\right)
                \end{aligned}
                $$
                On conclut que $c_n = 0$ en remarquant que $\omega_0 T = 2\pi$ et donc  $$e^{-\complexj (n-1)\omega_0 T}= e^{-\complexj (n+1)\omega_0 T} =1.$$

            \end{solution}
            \item \textit{(Pour les étudiants en avance.)} Calculer $c_1$ et $c_{-1}$. Indice : $\underset{x\to 0}{\lim} \cfrac{\sin(x)}{x} = 1$.
            \begin{solution}
                Pour $n=-1$,
                $\cfrac{1 - e^{-\complexj (n-1)\omega_0 T}}{\complexj (n-1)\omega_0} = 0$
                mais $\cfrac{1 - e^{-\complexj (n+1)\omega_0 T}}{\complexj (n+1)\omega_0}$ n'est pas définie.

                On remarquera toutefois que
                $$
                \begin{aligned}
                    \frac{e^{\complexj\varphi}}{2T}\cfrac{1 - e^{-\complexj (n+1)\omega_0 T}}{\complexj (n+1)\omega_0} &= e^{-\complexj ((n+1)\pi - \varphi)}\times \cfrac{e^{\complexj (n+1)\pi} - e^{-\complexj (n+1)\pi}}{2\complexj (n+1)\pi} \\
                    & = e^{-\complexj ((n+1)\pi - \varphi)}\times \frac{\sin((n+1)\pi)}{(n+1)\pi}
                \end{aligned}
                    $$

                La limite donnée en indice nous permet donc d'écrire par passage à la limite
                $$
                    c_{-1} = e^{j\varphi} 
                $$
                De la même façon on trouve
                $$
                    c_{1} = e^{j\varphi} 
                $$
            \end{solution}
        \end{enumerate}
        
\end{exercice}


%% ── Ex 2–3 : PROVE ───────────────────────────────────────────────
%% Purpose: make students re-derive a theorem stated in the CM.
%% Structure: easy direction → hard direction → verify on an example.






%% ── Ex 4–5 : APPLY ───────────────────────────────────────────────
%% Purpose: numerical problem with real or realistic GEII data.
%% Ends with an engineer's conclusion, not just a number.
\begin{exercice}[Impulsion périodique étroite --- approfondissement]
    Soit $0<\theta<T$ et $f_\theta$ la fonction $T$-périodique définie, sur la période $\left[-\dfrac T2,\dfrac T2\right]$, par
    $$
        f_\theta(t) = \begin{cases} \dfrac{1}{\theta} & \text{si } |t|\leq\dfrac{\theta}{2} \\ 0 & \text{si } \dfrac{\theta}{2} < |t| \leq \dfrac{T}{2} \end{cases},
    $$
    et prolongée à $\R$ tout entier par $T$-périodicité (un motif d'aire $1$ par période, concentré sur une largeur $\theta$).

    \begin{center}
    \begin{tikzpicture}[x=0.9cm,y=1.2cm]
        \draw[-Latex] (-1,0) -- (13,0) node[right] {$t$};
        \draw[-Latex] (0,-0.3) -- (0,3.3) node[above] {$f_\theta(t)$};
        \draw[niceblue,thick] (-1,0) -- (13,0);
        \foreach \c in {0,6,12} {
            \draw[niceblue,thick] (\c-0.5,0) -- (\c-0.5,2.5) -- (\c+0.5,2.5) -- (\c+0.5,0);
        }
        \draw[dashed] (5.5,0) -- (5.5,2.9);
        \draw[dashed] (6.5,0) -- (6.5,2.9);
        \draw[latex-latex] (5.5,2.9) -- (6.5,2.9) node[midway,above] {\small $\theta$};
        \node[left] at (0,2.5) {$\dfrac1\theta$};
        \node[below] at (6,0) {$T$};
        \node[below] at (12,0) {$2T$};
    \end{tikzpicture}
    \end{center}

    \begin{enumerate}[\bfseries 1.]
        \item Quelle est la période de $f_\theta$ ?
        \begin{solution}
            $T$, par construction.
        \end{solution}

        \item Calculer $a_0$.
        \begin{solution}
            Par définition, $a_0$ s'obtient en intégrant sur une période entière, ici $\left[-\frac T2,\frac T2\right]$. Par la relation de Chasles, on découpe cette période aux points où $f_\theta$ change d'expression :
            $$
                a_0 = \frac1T\int_{-T/2}^{T/2} f_\theta(t)\,\d t
                = \frac1T\left(\int_{-T/2}^{-\theta/2} 0\,\d t \ + \int_{-\theta/2}^{\theta/2} \frac{1}{\theta}\,\d t \ + \int_{\theta/2}^{T/2} 0\,\d t\right).
            $$
            Les deux intégrales aux extrémités sont nulles ($f_\theta=0$ sur ces régions) ; il ne reste que
            $$
                a_0 = \frac1T\int_{-\theta/2}^{\theta/2} \frac{1}{\theta}\,\d t = \frac{1}{T\theta}\times\theta = \frac1T.
            $$
        \end{solution}

        \item Calculer les coefficients de la série de Fourier $a_n$ et $b_n$, puis en déduire les valeurs du spectre $A_n$ en fonction de $\theta$.
        \begin{solution}
            $f_\theta$ est paire, donc $b_n=0$ pour tout $n$. Pour $a_n$, le même découpage qu'à la question précédente ne laisse subsister que l'intégrale sur $\left[-\frac\theta2,\frac\theta2\right]$ :
            $$
                a_n = \frac{2}{T}\int_{-\theta/2}^{\theta/2} \frac{1}{\theta}\cos(n\omega_0 t)\,\d t = \frac{2}{T\theta}\left[\frac{\sin(n\omega_0 t)}{n\omega_0}\right]_{-\theta/2}^{\theta/2} = \frac{2\sin\left(n\omega_0\theta/2\right)}{Tn\omega_0\theta/2}.
            $$
            En notant $\sinc(x) = \dfrac{\sin(x)}{x}$, et puisque $\omega_0\theta/2 = \pi\theta/T$,
            $$
                a_n = \frac{2}{T}\sinc\!\left(\frac{n\pi\theta}{T}\right), \qquad A_n = |a_n| = \frac{2}{T}\left|\sinc\!\left(\frac{n\pi\theta}{T}\right)\right|.
            $$
        \end{solution}

        \item Calculer les coefficients de la représentation complexe $c_n$ en fonction de $\theta$.
        \begin{solution}
            $b_n=0$ donc $c_n = a_n/2$ pour $n\geq1$, et $c_0=a_0$. En remarquant que $\sinc(0)=1$, on peut écrire une seule formule valable pour tout $n\in\Z$ :
            $$
                c_n = \frac1T\sinc\!\left(\frac{n\pi\theta}{T}\right).
            $$
        \end{solution}

        \item On fait maintenant tendre $\theta$ vers $0$, en gardant la hauteur $1/\theta$ (donc l'aire du motif reste égale à $1$). Que devient $c_n$ pour un $n$ fixé ?
        \begin{solution}
            Quand $\theta\to0$, $\dfrac{n\pi\theta}{T}\to0$ donc $\sinc\!\left(\dfrac{n\pi\theta}{T}\right)\to1$, et
            $$
                c_n \longrightarrow \frac1T \quad \text{pour tout } n\in\Z.
            $$
            Le spectre devient plat : toutes les fréquences $nf_0$ contribuent avec le même poids $1/T$, sans décroissance. C'est l'inverse de $f_{\rm carr\acute{e}}$, dont le spectre décroît en $1/n$ (énergie concentrée sur les basses fréquences) : une impulsion infiniment brève (mais d'aire finie) répartit son énergie sur toutes les fréquences à parts égales. On retrouvera ce résultat plus tard dans le cours.
        \end{solution}
    \end{enumerate}
\end{exercice}


\begin{exercice}[Série de Fourier complexe d'un signal réel]

    Soit $f$ une fonction $T$-périodique à valeurs réelles, vérifiant les conditions de Dirichlet, de coefficients de Fourier complexes $c_n$. On veut montrer que, bien que les $c_n$ soient a priori complexes, la série $\displaystyle\sum_{n=-\infty}^{+\infty} c_n e^{jn\omega_0 t}$ est réelle pour tout $t$.

    \begin{enumerate}[\bfseries 1.]
        \item Pour $n\in\N$, écrire l'expression intégrale de $c_{-n}$.
        \begin{solution}
            Il suffit de remplacer $n$ par $-n$ dans la définition :
            $$
                c_{-n} = \frac1T\int_0^T f(t)e^{jn\omega_0t}\,\d t.
            $$
        \end{solution}

        \item Calculer $\overline{c_n}$ en utilisant le fait que $f$ est à valeurs réelles ($\overline{f(t)}=f(t)$) et conclure que $c_{-n}=\overline{c_n}$.
        \begin{solution}
            $$
                \overline{c_n} = \overline{\frac1T\int_0^T f(t)e^{-jn\omega_0t}\,\d t} = \frac1T\int_0^T \overline{f(t)}\,\overline{e^{-jn\omega_0t}}\,\d t = \frac1T\int_0^T f(t)e^{jn\omega_0t}\,\d t.
            $$
            On reconnaît exactement l'expression de $c_{-n}$ trouvée à la question précédente : $\overline{c_n} = c_{-n}$.
        \end{solution}

        \item Montrer que $c_n e^{jn\omega_0t} + c_{-n}e^{-jn\omega_0t}$ est un nombre réel.
        \par \textbf{Rappel :} $\overline{e^{\complexj t}} = e^{-\complexj t}$ et pour tout complexe $z\in\C$, $z + \overline{z} = 2\,\mathrm{Re}(z) \in \R$.
        \begin{solution}
            D'après la question précédente, en posant $z=c_ne^{jn\omega_0t}$,
            $$
                c_n e^{jn\omega_0t} + c_{-n}e^{-jn\omega_0t} = z + \overline{z}.
            $$
            Or pour tout $z\in\C$, $z+\overline{z}=2\,\mathrm{Re}(z)\in\R$. La somme est donc réelle.
        \end{solution}

        \item Conclure que $\displaystyle\sum_{n=-\infty}^{+\infty} c_n e^{jn\omega_0 t} \in \R$ pour tout $t\in\R$.
        \begin{solution}
            On regroupe les termes d'indices $n$ et $-n$ ($n\geq1$) :
            $$
                \sum_{n=-\infty}^{+\infty} c_n e^{jn\omega_0 t} = c_0 + \sum_{n=1}^{+\infty}\left(c_n e^{jn\omega_0t} + c_{-n}e^{-jn\omega_0t}\right).
            $$
            Le terme $c_0=\frac1T\int_0^T f(t)\,\d t$ est réel ($f$ à valeurs réelles), et chaque terme groupé est réel (question précédente) : la série est une somme de nombres réels, donc réelle.
        \end{solution}

    \end{enumerate}

    On applique maintenant ce résultat pour démontrer le Corollaire~3.1 du cours : si $f$ est paire, $b_n=0$ pour tout $n$ (série de cosinus) ; si $f$ est impaire, $a_0=0$ et $a_n=0$ pour tout $n$ (série de sinus).

    \begin{enumerate}[\bfseries 1.]
        \setcounter{enumi}{4}
        \item Justifier que l'on peut calculer $a_0$, $a_n$ et $b_n$ en intégrant sur $[-\frac{T}{2},\frac{T}{2}]$ plutôt que sur $[0,T]$.
        \begin{solution}
            L'intégrale d'une fonction $T$-périodique ne dépend pas du choix de l'intervalle d'intégration, pourvu qu'il soit de longueur $T$. L'intervalle $[-\frac{T}{2},\frac{T}{2}]$ est de longueur $T$, donc convient.
        \end{solution}

        \item On suppose $f$ paire. Montrer que $t\mapsto f(t)\sin(n\omega_0 t)$ est impaire, puis en déduire que $b_n=0$ et que $c_n$ est réel pour tout $n$.
        \begin{solution}
            $f(-t)\sin(-n\omega_0 t) = f(t)\times\left(-\sin(n\omega_0 t)\right) = -f(t)\sin(n\omega_0 t)$ : produit d'une fonction paire ($f$) par une fonction impaire ($\sin$), donc impair. D'après la Propriété~1.3 et la question précédente,
            $$
                b_n = \frac{2}{T}\int_{-T/2}^{T/2} f(t)\sin(n\omega_0 t)\,\d t = 0.
            $$
            Or $c_n = \dfrac{a_n-\complexj b_n}{2}$ (cours) ; avec $b_n=0$, $c_n=\dfrac{a_n}{2}$, qui est réel puisque $a_n$ est une intégrale d'une fonction réelle.
        \end{solution}

        \item On suppose maintenant $f$ impaire. Montrer que $a_0=0$, puis que $t\mapsto f(t)\cos(n\omega_0 t)$ est impaire pour en déduire que $a_n=0$ et que $c_n$ est imaginaire pur pour tout $n\geq1$.
        \begin{solution}
            $a_0=\dfrac1T\displaystyle\int_{-T/2}^{T/2} f(t)\,\d t = 0$ directement par la Propriété~1.3 ($f$ impaire sur un intervalle symétrique). Ensuite $f(-t)\cos(-n\omega_0 t) = -f(t)\cos(n\omega_0 t)$ : produit d'une fonction impaire ($f$) par une fonction paire ($\cos$), donc impair. D'après la Propriété~1.3,
            $$
                a_n = \frac{2}{T}\int_{-T/2}^{T/2} f(t)\cos(n\omega_0 t)\,\d t = 0.
            $$
            Or $c_n = \dfrac{a_n-\complexj b_n}{2}$ (cours) ; avec $a_n=0$, $c_n=-\dfrac{\complexj b_n}{2}$, qui est imaginaire pur puisque $b_n$ est réel.
        \end{solution}
    \end{enumerate}
\end{exercice}






%% ── Ex 6–7 (optional) : EXTEND ───────────────────────────────────
%% Purpose: harder result not in the CM, for strong students.
%% Mark clearly as optional difficulty.
%
% \begin{exercice}[Approfondissement]
%     \begin{enumerate}[\bfseries 1.]
%         \item ...
%         \begin{solution}
%         \end{solution}
%     \end{enumerate}
% \end{exercice}

\end{document}
